O algoritmo Aho-Corasick é um pilar dos sistemas de detecção de intrusão de
rede baseados em assinaturas e das ferramentas de análise de malware, pois
varre uma entrada em tempo linear enquanto casa muitos padrões
simultaneamente. Sua travessia canônica do autômato, porém, é sequencial e
subutiliza os processadores multinúcleo modernos quando uma única entrada de
grande porte precisa ser inspecionada. Este trabalho projeta, implementa e
avalia estratégias de varredura Aho-Corasick com paralelismo de dados,
construídas sobre POSIX Threads para CPUs multinúcleo de memória
compartilhada. O texto de entrada é particionado em segmentos sobrepostos que
as threads varrem concorrentemente sobre um autômato compartilhado e
estritamente somente-leitura, mantendo a busca livre de travas; propõe-se
também um layout de saída achatado e contíguo que reduz o custo de
perseguição de ponteiros na emissão de casamentos. A avaliação foca o regime
limitado por memória, no qual conjuntos de regras realistas do Snort e do
Emerging Threats produzem autômatos de até cerca de $507$~MiB, muito além da    
capacidade efetiva da cache de último nível. Em uma estação de trabalho
homogênea AMD Ryzen~9~9950X (Zen~5, dezesseis núcleos, trinta e duas threads
de hardware), a estratégia dinâmica com saída achatada atinge $22{,}91\times$
de speedup e $7\,542$~MB/s no dicionário completo do Snort, e $18{,}96\times$
e $3\,979$~MB/s no do Emerging Threats. A mesma estratégia permanece
insensível ao conteúdo do corpus na carga do Snort, com $7\,976$~MB/s no
SimpleWiki e $7\,595$~MB/s no Enron, e a granularidade de tarefa mais grossa
testada atinge $7\,804$~MB/s. De ponta a ponta, sobre o corpus Enron replicado
oito vezes, as acelerações atingem $22{,}23\times$ para o Snort e
$17{,}71\times$ para o Emerging Threats, mostrando que a construção permanece
um custo limitado e que a busca domina o pipeline. Um diagnóstico controlado
sobre corpora com conteúdo enviesado --- mesmos bytes e mesmo total de
casamentos, concentrados em uma região do texto --- confirma o balanceamento
de carga da estratégia dinâmica: partições estáticas perdem até cerca de
$30\%$
as ultrapassa. A conclusão central é que, quando o autômato excede a
capacidade efetiva da cache, a escalabilidade do Aho-Corasick passa a ser
limitada pelo subsistema de memória, e não pela lógica de casamento; o
paralelismo de texto sobre um autômato compartilhado, combinado à emissão
achatada de casamentos, é a resposta mais eficaz nesse regime.
